%\ifodd\value{page}\hbox{}\newpage\fi
\chapter*{Ṛgveda 10.125 — Devī Sūkta, Mantra 5}
\addcontentsline{toc}{chapter}{Devī Sūkta 10.125.5}

\begin{sloka}
{\sanskritfont
\begin{tabular}{c}
अहं रुद्राय धनुरातनोमि ब्रह्मद्विषे शरवे हन्तवा उ ।\\
अहं जनाय समदं कृणोम्यहं द्यावापृथिवी आविवेश ॥५॥
\end{tabular}

\vspace{1.5em}

{\middlesize \iastfont
ahaṁ rudrāya dhanur ātanomi brahmadviṣe śarave hantavā u |\\
ahaṁ janāya samadaṁ kṛṇomy ahaṁ dyāvāpṛthivī āviveśa ||5||}

\vspace{1em}

{\middlesize
\anvaya{%
अहं रुद्राय धनुः आतनोमि ।
ब्रह्मद्विषे शरवे हन्तवाः उ ।
अहं जनाय समदम् कृणोमि ।
अहं द्यावापृथिवी आविवेश ।
}
}

\middlesize \iastfont \itmean{%
«Io tendo l’arco per Rudra,  
affinché la freccia colpisca chi odia il sacro.  
Io preparo il conflitto per gli uomini;  
io pervado il Cielo e la Terra.»%
}

}
\end{sloka}

\wordmeaning{%
\wordline{अहं रुद्राय धनुः आतनोमि}
{ahaṁ rudrāya dhanuḥ ātanomi}
{ io tendo (\textit{ātanomi}) l’arco (\textit{dhanuḥ}) per Rudra, il principio della forza trasformante; }%

\wordline{ब्रह्मद्विषे शरवे हन्तवाः}
{brahmadviṣe śarave hantavāḥ}
{ affinché la freccia (\textit{śara}) colpisca (\textit{hantavāḥ}) chi odia (\textit{dviṣ}) il \textit{brahman}, il sacro ordine e la parola rivelata; }%

\wordline{अहं जनाय समदम् कृणोमि}
{ahaṁ janāya samadam kṛṇomi}
{ io preparo (\textit{kṛṇomi}) il confronto, il conflitto (\textit{samada}) per gli uomini (\textit{jana}); }%

\wordline{अहं द्यावापृथिवी आविवेश}
{ahaṁ dyāvāpṛthivī āviveśa}
{ io pervado (\textit{āviveśa}) il Cielo e la Terra (\textit{dyāvā-pṛthivī}), l’intera dimensione cosmica; }%
}

Nel quinto mantra la Dea manifesta il suo volto più deciso e attivo. Lei non è soltanto parola, conoscenza o sovranità silenziosa, ma anche forza che prende parte al conflitto. Tendere l’arco per \textit{Rudra} significa rendere efficace l’energia distruttiva che elimina ciò che si oppone al \textit{brahman}, inteso come ordine sacro e verità cosmica. La preparazione del \textit{samada}, il confronto o scontro, non indica una violenza cieca, ma la tensione necessaria a ristabilire equilibrio. Il conflitto è presentato come parte integrante del processo cosmico, non come sua negazione. Il mantra si chiude con un’affermazione di pervasività totale: la Dea “entra” e permea Cielo e Terra. La forza che combatte e quella che sostiene non sono due principi distinti. È la stessa potenza che agisce nel conflitto e che abita l’intero cosmo, rendendo possibile sia la distruzione del disordine sia la continuità dell’esistenza.

Il termine \textit{samada} introduce nel \textit{Devī Sūkta} una concezione del conflitto profondamente diversa da quella di una violenza cieca o distruttiva. \textit{Samada} indica un “andare incontro”, uno scontro frontale tra forze contrapposte. Quando la Dea afferma \textit{ahaṁ janāya samadaṁ kṛṇomi}, non dichiara di seminare caos, bensì di preparare la tensione necessaria al ristabilimento dell’equilibrio. Il conflitto non nasce dall’arbitrio, ma dall’esistenza di ciò che si oppone al \textit{brahman}, all’ordine sacro e alla verità che sostiene il mondo. In questo senso, \textit{samada} è una funzione discriminante: separa, espone, mette alla prova. È il momento in cui il disordine viene portato alla luce per poter essere rimosso. Il fatto che lo stesso soggetto che prepara il \textit{samada} dichiari subito dopo di pervadere Cielo e Terra mostra che distruzione e sostegno non sono principi opposti, ma aspetti di una medesima potenza. \textit{Samada} nomina così il volto severo della compassione cosmica: la forza che non distrugge per annientare, ma per rendere possibile la continuità dell’ordine e della vita.
